\section*{Chapitre \ref{ch:g8:the-oceans} --- Les océans}

\begin{solution}{exo:g8:the-oceans:1}
\omterm{def:g8:the-oceans:seafloor}{Plateau continental}, \omterm{def:g8:the-oceans:seafloor}{talus continental}, \omterm{def:g8:the-oceans:seafloor}{plaine abyssale}, \omterm{def:g8:the-oceans:seafloor}{dorsale océanique}.
\end{solution}

\begin{solution}{exo:g8:the-oceans:2}
(a) Faux~: environ \qty{3700}{m}, soit \qty{3.7}{km}. (b) Vrai. (c)
Faux~: l'évaporation enlève de l'eau et laisse les sels~: l'eau devient
donc plus salée. (d) Vrai.
\end{solution}

\begin{solution}{exo:g8:the-oceans:3}
(a) Une \omterm{def:g8:the-oceans:seafloor}{dorsale océanique}. (b) La \omterm{def:g8:the-oceans:salinity}{salinité}. (c) Un \omterm{def:g8:the-oceans:current}{gyre}. (d) Une \omterm{def:g8:the-oceans:seafloor}{fosse océanique}.
\end{solution}

\begin{solution}{exo:g8:the-oceans:4}
(a) Dans le sens inverse des aiguilles d'une montre. (b) Chaud~: le Gulf
Stream~; froid~: le courant des Canaries. (c) Les \omterm{def:g8:winds-global-circulation:global-winds}{alizés}.
\end{solution}

\begin{solution}{exo:g8:the-oceans:5}
(a) $35 \div 1000 \times 100 = 3.5~\%$. (b) $20 \times 35 = \qty{700}{g}$.
(c) $1000 \div 35 \approx 28.6$~: environ \qty{29}{kg}.
\end{solution}

\begin{solution}{exo:g8:the-oceans:6}
(a) $1500 \times 6 = \qty{9000}{m}$. (b) Le son descend jusqu'au fond et
remonte~: il parcourt deux fois la profondeur. (c) $9000 \div 2 =
\qty{4500}{m}$. (d) Une \omterm{def:g8:the-oceans:seafloor}{plaine abyssale} (une fosse serait plus profonde,
une dorsale ou un plateau moins profonds).
\end{solution}

\begin{solution}{exo:g8:the-oceans:7}
(a) Entre environ \qty{3500}{m} et \qty{5000}{m}. (b) À environ \qty{4400}{km} le long du profil, à environ \qty{2300}{m} de profondeur. (c) Environ
$4500 - 2300 = \qty{2200}{m}$ (entre 1200 et \qty{2700}{m}, selon la
\omterm{def:g1:landscapes-around-us:landform}{plaine} choisie). (d) L'\omterm{def:g2:water-on-earth:ocean}{océan} a quelques kilomètres de profondeur mais des
milliers de kilomètres de largeur~: dessiné à la même échelle, tout le
profil serait plus mince qu'un trait.
\end{solution}

\begin{solution}{exo:g8:the-oceans:8}
La \omterm{def:g2:water-on-earth:ocean}{mer} A est plus salée~: le Soleil et l'air sec provoquent une forte
\emph{évaporation}, l'eau s'en va, \emph{les sels restent}, et presque
aucune \omterm{def:g4:clouds-rain-snow:precipitation}{pluie} ni aucun fleuve ne rapporte d'\emph{\omterm{def:g2:water-on-earth:salt-fresh}{eau douce}}. La \omterm{def:g2:water-on-earth:ocean}{mer} B est
moins salée~: ses fleuves et sa \omterm{def:g4:clouds-rain-snow:precipitation}{pluie} apportent beaucoup d'\emph{\omterm{def:g2:water-on-earth:salt-fresh}{eau douce}},
qui dilue les sels.
\end{solution}

\begin{solution}{exo:g8:the-oceans:9}
(a) Le Gulf Stream, vers le nord le long de la \omterm{def:g1:landscapes-around-us:water}{côte}, puis vers le
nord-est. (b) Vers l'est (le nord-est), avec la dérive nord-atlantique.
(c) Le long de la \omterm{def:g1:landscapes-around-us:water}{côte} du nord-ouest de l'Afrique (et de l'Espagne et du
Portugal), avec le courant des Canaries. (d) Froid.
\end{solution}

\begin{solution}{exo:g8:the-oceans:10}
(a) $6.4 \times 24 = 153.6$~: environ \qty{154}{km}. (b) $153.6 \times 7
\approx \qty{1075}{km}$. (c) Le Gulf Stream est un peu plus rapide qu'un
marcheur, mais il ne s'arrête jamais, ni le jour ni la nuit.
\end{solution}

\begin{solution}{exo:g8:the-oceans:11}
(a) \qty{3.7}{km} et \qty{11}{km}. (b) $11000 \div 3700 \approx 3.0$
fois. (c) Le son parcourt $2 \times 11000 = \qty{22000}{m}$~; $22000
\div 1500 \approx 14.7$~: environ \qty{15}{s}. (d) $1000 = 10^3$~ans.
\end{solution}

\begin{solution}{exo:g8:the-oceans:12}
(a) Près de l'\omterm{def:g3:continents-and-oceans:equator}{équateur}, les \omterm{def:g8:winds-global-circulation:global-winds}{alizés} poussent l'eau tropicale chaude vers
l'ouest. Au bord ouest de l'Atlantique, elle remonte vers le nord avec le
Gulf Stream, un courant chaud, le long de la Floride et de l'Amérique du
Nord. Les \omterm{def:g8:winds-global-circulation:global-winds}{vents d'ouest} la poussent ensuite à travers l'Atlantique vers le
nord-est, jusqu'aux \omterm{def:g2:water-on-earth:ocean}{mers} situées au large de la Norvège. (b) Là, elle cède
sa chaleur à l'air froid~: elle se refroidit, et elle devient plus salée
quand la banquise se forme, parce que le sel reste dans l'eau. L'eau
froide et salée est dense~: elle coule donc jusqu'au fond. (c) Elle se
glisse vers le sud le long du fond de l'\omterm{def:g2:water-on-earth:ocean}{océan}, vers l'Antarctique, puis
dans les \omterm{def:g2:water-on-earth:ocean}{océans} Indien et Pacifique, où elle remonte lentement.
\end{solution}

\begin{solution}{pb:g8:the-oceans:1}
\textbf{1.} Dans l'\omterm{def:g2:water-on-earth:ocean}{océan} Atlantique, dans l'\omterm{def:g3:continents-and-oceans:equator}{hémisphère} Nord.

\textbf{2.} Dans le sens des aiguilles d'une montre~; la bouée est d'abord emportée vers le nord.

\textbf{3.} L'étape A (Gulf Stream) est chaude~: elle vient des \omterm{def:g5:weather-and-climate:tropics}{tropiques}.
L'étape C (courant des Canaries) est froide~: elle vient de plus hautes
\omterm{def:g6:latitude-longitude:latitude}{latitudes} et va vers l'\omterm{def:g3:continents-and-oceans:equator}{équateur}.

\textbf{4.} Étape B~: les \omterm{def:g8:winds-global-circulation:global-winds}{vents d'ouest}. Étape D~: les \omterm{def:g8:winds-global-circulation:global-winds}{alizés}.

\textbf{5.} $3800 \div 100 = 38$~jours.

\textbf{6.} B~: $2100 \div 20 = 105$~jours. C~: $3100 \div 20 = 155$~jours.
D~: $6100 \div 20 = 305$~jours.

\textbf{7.} $3800 + 2100 + 3100 + 6100 = \qty{15100}{km}$.

\textbf{8.} Le Gulf Stream est le courant étroit et rapide du côté ouest du
\omterm{def:g8:the-oceans:current}{gyre}~; les courants des autres côtés sont larges et lents.

\textbf{9.} $38 + 105 + 155 + 305 = 603$~jours.

\textbf{10.} $603 \div 30 = 20.1$~: environ \textbf{20~mois}.

\textbf{11.} Après les étapes A, B et C, $38 + 105 + 155 = 298$~jours se
sont écoulés~; le 365\textsuperscript{e}~jour arrive pendant l'étape D (le
courant nord-équatorial), $365 -
298 = 67$~jours après son début.

\textbf{12.} $603 - 365 = 238$~jours de plus, soit $238 \div 30 \approx 8$
mois de plus.
\end{solution}
